\section{Conclusões}

\begin{itemize}
    \item Para o Cenário 25, linha de montagem de ar-condicionado com 12 postos, o sistema desenvolvido pela Equipe 4 integra o CLP LS XGK, a Raspberry Pi e a supervisão no Windows por protocolos padronizados (Modbus RTU, OPC UA com \cod{Basic256Sha256}/\cod{SignAndEncrypt}, SNTP) e por um protocolo binário de Andon de 16 bytes com ACK. O Desktop fala os mesmos protocolos do Gateway, de modo que o receptor e a Dashboard funcionam sem alteração de código com a linha real ou com o simulador. No simulador, porém, não há os nós e métodos \cod{Linha/*}: o estado dos postos, o contador de produção, a IT Digital sincronizada e os comandos de base, Start e reset exigem o Gateway.
    \item Na execução com o CLP real, os 18 eventos de Andon foram capturados, confirmados ao CLP em 13,4 a 17,3 ms e entregues ao receptor em 23,4 a 39,9 ms, sem nova tentativa. Durante uma perda de comunicação de 1 h 38 min, quatro eventos ficaram retidos na FIFO do CLP e foram entregues em ordem na reconexão: a combinação ``fila no Ladder + COMMIT antes do \cod{P10000} + fila SQLite no Gateway + ACK TCP após persistência'' não perdeu nem duplicou eventos.
    \item O Gateway mostrou-se resiliente: manteve OPC UA e TCP ativos sem o CLP, tentou reabrir a porta a cada 2 s (4.444 tentativas registradas), sincronizou o relógio do CLP logo após a reconexão e encerrou de forma ordenada.
    \item A Dashboard calcula Disponibilidade, Desempenho, Qualidade e OEE a partir dos estados e do contador do CLP; a conferência manual da Figura \ref{fig:dash3} reproduz os valores exibidos (84,1 \%, 40,0 \%, 66,7 \% e 22,4 \%). Os estados dos 12 postos tornam visíveis o fluxo de peças e o gargalo (P07), objetivo central do acompanhamento \textit{Lean}.
    \item A IT Digital E07 acompanhou o ciclo do P07 no CLP emulado com trocas de slide até 0,9 s após o limite teórico e respeitou as regras de parada por Andon e de espera pelo próximo ciclo.
    \item A API rejeitou comandos de outra origem (403) e dados fora do domínio (422), e os comandos de base, Start e reset só foram confirmados após a verificação no CLP. As 162 verificações automatizadas aprovadas (98 do Simulador e 64 do Gateway) cobrem os casos de falha difíceis de reproduzir no hardware.
    \item Limitações: as capturas de tela e as medições da Dashboard e da IT Digital foram obtidas com o Gateway real e um CLP emulado, que simplifica o Ladder; as redes completas dos postos 03 a 06 e 12 não constam nas capturas do XG5000; e o protocolo de Andon não tem identificador de evento, o que impede garantir processamento exatamente uma vez. O gêmeo digital e a conexão entre gêmeos (T2T) não foram implementados, e o log e os bancos originais da execução real não acompanham a entrega, que inclui apenas o extrato dos 18 eventos. A amostra de ciclos físicos da E07 ainda não foi coletada, portanto não há comparação estatística entre os ciclos reais e os simulados, e a IT Digital, de 36 s, excede a meta de 34 s da E07 e o Takt de 34,6 s do cenário. Recomenda-se repetir as capturas com a linha física completa, coletar ao menos 30 ciclos físicos, compatibilizar a duração da IT e incluir um identificador no protocolo antes de uma homologação.
\end{itemize}
